\section{Background: Farid's Algorithms 1 and 2}
\label{sec:background}

This section restates the two algorithms as printed in
Farid et al.~\cite{farid2014hybrid} (Section 4, pp.~1941--1942), because
the replication in Section~\ref{sec:setup} depends on reproducing their
exact steps rather than a paraphrase of their intent.

\subsection{Algorithm 1: the hybrid decision tree}

Algorithm 1 first classifies every training instance with a NB classifier
fit on the same training data, using the standard class-conditional
probabilities $P(C_i)$ and $P(A_{ij}\mid C_i)$; every instance the NB
classifier misclassifies (its posterior-probability class disagrees with
the recorded label) is removed from the training set. A C4.5-style
decision tree is then grown on the reduced training set by ordinary
recursive splitting: pick the best splitting attribute, create a labeled
root, branch on each split predicate, and recurse until a stopping
criterion is met at each branch.

\begin{algorithm}[t]
\caption{Hybrid decision tree induction (Farid et al.~\cite{farid2014hybrid}, Alg.\ 1)}
\label{alg:1}
\begin{algorithmic}[1]
\Require Training dataset $D=\{x_1,\dots,x_n\}$ with class labels
\Ensure Decision tree $T$
\For{each class $C_i \in D$}
  \State find prior probability $P(C_i)$
\EndFor
\For{each attribute value $A_{ij} \in D$}
  \State find class-conditional probability $P(A_{ij}\mid C_i)$
\EndFor
\For{each training instance $x_i \in D$}
  \State find posterior probability $P(C_i\mid x_i)$
  \If{$x_i$ is misclassified}
    \State remove $x_i$ from $D$
  \EndIf
\EndFor
\State $T \gets \emptyset$; determine the best splitting attribute
\State create the root node of $T$, labeled with that attribute
\State add an arc to the root for each split predicate/label
\For{each arc}
  \State $D \gets$ the subset of $D$ that satisfies the arc's predicate
  \If{a stopping point is reached on this path}
    \State create a leaf labeled with the majority class
  \Else
    \State recursively build a subtree on $D$
  \EndIf
  \State attach the resulting subtree to the arc
\EndFor
\end{algorithmic}
\end{algorithm}

\subsection{Algorithm 2: the hybrid Na\"ive Bayes classifier}

Algorithm 2 reverses the order: a decision tree is grown first on the full
training set, by the same recursive splitting procedure as Algorithm~1's
second half. Every attribute the tree tests is then given a weight
$W_i = 1/\sqrt{d}$, where $d$ is the smallest depth (root $=1$) at which
the tree tests that attribute; an attribute the tree never tests gets
$W_i = 0$ and is dropped from the classifier entirely. A NB classifier is
then built using only the attributes with $W_i \neq 0$, with each
class-conditional probability raised to the power of its attribute's
weight:
\begin{equation}
P(x_i \mid C_i) = \prod_{j=1}^{n} P(A_{ij}\mid C_i)^{W_i}. \label{eq:14}
\end{equation}
Equation~\eqref{eq:14} is the paper's Eq.~(14); the posterior
$P(C_i\mid x_i)$ used for classification combines this weighted likelihood
with the class prior in the usual NB way.

\begin{algorithm}[t]
\caption{Hybrid Na\"ive Bayes classifier (Farid et al.~\cite{farid2014hybrid}, Alg.\ 2)}
\label{alg:2}
\begin{algorithmic}[1]
\Require Training dataset $D=\{x_1,\dots,x_n\}$ with class labels
\Ensure A NB classification model
\State build a decision tree $T$ on $D$ (as in Alg.~\ref{alg:1}, lines 9--16)
\For{each attribute $A_i \in D$}
  \If{$A_i$ is not tested in $T$}
    \State $W_i \gets 0$
  \Else
    \State $d \gets$ the smallest depth at which $T$ tests $A_i$
    \State $W_i \gets 1/\sqrt{d}$
  \EndIf
\EndFor
\For{each class $C_i \in D$}
  \State find prior probability $P(C_i)$
\EndFor
\For{each attribute $A_i \in D$ with $W_i \neq 0$}
  \For{each value $A_{ij}$ of $A_i$}
    \State find $P(A_{ij}\mid C_i)$
  \EndFor
\EndFor
\For{each instance $x_i \in D$}
  \State find posterior probability $P(C_i\mid x_i)$ using Eq.~\eqref{eq:14}
\EndFor
\end{algorithmic}
\end{algorithm}

Two properties of these algorithms matter for what follows. First, neither
step in Algorithm~1 or Algorithm~2 is stated as being confined to a
training fold; the paper's description is silent on whether the filter or
the tree ever sees an instance that is later scored as a test case. Second,
Algorithm 1's notion of noise is entirely internal to one NB fit: an
instance is noise if and only if this particular classifier, trained on
this particular data, disagrees with its label, with no second opinion and
no confidence threshold.
